% !TEX root = ../../../main.tex
% 所属：第四部分：Fourier 分析 / 第1章：Fourier级数的基本性质
% 节：1.2 Fourier级数的唯一性
% 正文按原文件顺序保存；由 ../chapter1.tex 调用。
\section{Fourier级数的唯一性}
\begin{theorem}\label{Fourier级数的唯一性}
	假设 $f$ 是在圆周上可积的函数,且对所有 $n\in\Z$,有
	\[
	\widehat f(n)=0.
	\]
	若 $f$ 在点 $\theta_0$ 连续,则有
	\[
	f(\theta_0)=0.
	\]
\end{theorem}

\begin{proof}
	先假设 $f$ 是实值的。由于问题是周期的,并且可以通过平移把连续点移到原点,所以只需证明 $\theta_0=0$ 的情形。使用反证法,若 $f(0)\neq0$,不是一般性的可设$f(0)>0.$

	核心的思路是构造一个在零点取最大值的三角多项式列 $\{p_k\}$,使得当 $k\to\infty$ 时
	\[
	\int_{-\pi}^{\pi}p_k(\theta)f(\theta)\dd\theta\to\infty.
	\]
	但是根据 Fourier 系数全为零的假设,这些积分又都等于零,由此得到矛盾。
	
	由于 $f$ 在 $0$ 处连续,故可以选择 $0<\delta\leq \pi/2$,使得当 $|\theta|<\delta$ 时,
	\[
	f(\theta)>\frac{f(0)}2.
	\]
	令
	\[
	p(\theta)=\eps+\cos\theta,
	\]
	其中选择 $\eps>0$ 足够小,使当 $\delta\leq |\theta|\leq\pi$ 时,有
	\[
	|p(\theta)|<1-\frac\eps2.
	\]
	然后选择一个正数 $\eta$,且 $\eta<\delta$,使得对 $|\theta|<\eta$,有
	\[
	p(\theta)>1+\frac\eps2.
	\]
	最后令
	\[
	p_k(\theta)=\bigl[p(\theta)\bigr]^k.
	\]
	由$f\in \mathcal{R}(I)$,从而$f $有界,存在 $B>0$,使得对任意的 $\theta$,有 $|f(\theta)|\leq B$。每个 $p_k$ 都是三角多项式。由于对任意 $n$,有 $\widehat f(n)=0$,于是对任意 $k$,有
	\[
	\int_{-\pi}^{\pi}f(\theta)p_k(\theta)\dd\theta=0.
	\]
	另一方面,有估计
	\[
	\left|\int_{\delta\leq |\theta|\leq\pi}f(\theta)p_k(\theta)\dd\theta\right|
	\leq 2\pi B\left(1-\frac\eps2\right)^k.
	\]
	同时,选取 $\delta$ 使得当 $|\theta|<\delta$ 时,$p(\theta)$ 和 $f(\theta)$ 是非负的,于是
	\[
	\int_{\eta\leq |\theta|<\delta}f(\theta)p_k(\theta)\dd\theta\geq0.
	\]
	最后,在 $|\theta|<\eta$ 上,
	\[
	\begin{aligned}
		\int_{|\theta|<\eta}f(\theta)p_k(\theta)\dd\theta
		&\geq 2\eta\cdot \frac{f(0)}2\left(1+\frac\eps2\right)^k.
	\end{aligned}
	\]
	因此,当 $k\to\infty$ 时,
	\[
	\int_{-\pi}^{\pi}p_k(\theta)f(\theta)\dd\theta\to\infty,
	\]
	这与该积分恒等于零矛盾。所以实值情形下 $f(0)=0$。
	\begin{center}
		\begin{tikzpicture}[scale=1.0,>=Stealth]
			\draw[->] (-3.4,0) -- (3.4,0);
			\draw[->] (0,-0.2) -- (0,3.3);
			\draw[domain=-3.14:3.14,samples=200,smooth] plot (\x,{0.9+0.45*cos(\x r)});
			\draw[domain=-3.14:3.14,samples=300,smooth] plot (\x,{1.25*(0.1+cos(\x r))^3});
			\draw[domain=-1.2:1.2,samples=300,smooth] plot (\x,{2.45*(0.1+cos(\x r))^6});
			\node at (2.5,1.05) {$p$};
			\node at (0.95,1.9) {$p_3$};
			\node at (0.35,2.95) {$p_{6}$};
		\end{tikzpicture}
	\end{center}
	
	一般地,令
	\[
	f(\theta)=u(\theta)+iv(\theta),
	\]
	其中 $u$ 和 $v$ 是实值函数。若定义 $\overline f(\theta)=\overline{f(\theta)}$,那么有
	\[
	u(\theta)=\frac{f(\theta)+\overline f(\theta)}2,
	\qquad
	v(\theta)=\frac{f(\theta)-\overline f(\theta)}{2i}.
	\]
	又因为$\widehat{\overline f}(n)=\overline{\widehat f(-n)},$

	所以 $u$ 和 $v$ 的 Fourier 系数也全为零。由实值情形可得,在 $f$ 的连续点有
	\[
	u(\theta_0)=0,
	\qquad
	v(\theta_0)=0.
	\]
	因此 $f(\theta_0)=0$。
	
\end{proof}

根据上面的命题不难得到以下的命题
\begin{corollary}
	如果 $f$ 在圆周上连续且对所有 $n\in\Z$,有$\widehat f(n)=0,$则$f=0.$
	
\end{corollary}
\begin{corollary}
	如果 $f,g$ 在圆周上连续且对所有 $n\in\Z$,有$\widehat f(n)=\widehat g(n),$则$f=g.$
	
\end{corollary}
这也就是所谓的Fourier级数的唯一性

\begin{corollary}
	假设 $f$ 是在圆周上的连续函数且 $f$ 的 Fourier 系数绝对收敛,即
	\[
	\sum_{n=-\infty}^{\infty}|\widehat f(n)|<\infty.
	\]
	那么,$f$ 的 Fourier 级数一致收敛于 $f$,即
	\[
	\lim_{N\to\infty}S_N(f)(\theta)=f(\theta),
	\]
	且上述收敛关于 $\theta$ 是一致的。
\end{corollary}

\begin{proof}
由
	\[
	\sum_{n=-\infty}^{\infty}|\widehat f(n)|<\infty
	\]
	可知,$f$ 的 Fourier 系数的部分和绝对且一致收敛。事实上,由 Weierstrass $M$ 判别法,级数
	\[
	\sum_{n=-\infty}^{\infty}|\widehat f(n)e^{in\theta}| =\sum_{n=-\infty}^{\infty}|\widehat f(n)|<\infty
	\]
	在圆周上一致收敛。于是可以定义
	\[
	g(\theta)=\sum_{n=-\infty}^{\infty}\widehat f(n)e^{in\theta}
	=\lim_{N\to\infty}\sum_{n=-N}^{N}\widehat f(n)e^{in\theta}.
	\]
	由一致收敛和每一项的连续性可知,$g$ 在圆周上连续。
	
	进一步,由一致收敛可以交换无穷求和与积分,所以 $g$ 的 Fourier 系数恰好是 $\widehat f(n)$。具体地,
	\[
	\begin{aligned}
		\widehat g(n)
		&=\frac1{2\pi}\int_{-\pi}^{\pi}g(\theta)e^{-in\theta}\dd\theta                         \\
		&=\frac1{2\pi}\int_{-\pi}^{\pi}\sum_{m=-\infty}^{\infty}
		\widehat f(m)e^{im\theta}e^{-in\theta}\dd\theta                                             \\
		&=\frac1{2\pi}\sum_{m=-\infty}^{\infty}\widehat f(m)
		\int_{-\pi}^{\pi}e^{i(m-n)\theta}\dd\theta                                                   \\
		&=\widehat f(n).
	\end{aligned}
	\]
	因此 $f-g$ 的所有 Fourier 系数均为零。由上一推论,$f-g=0$,于是 $f=g$。也就是说,$f$ 的 Fourier 级数一致收敛于 $f$。
\end{proof}

\begin{corollary}\label{f(n)=O(1/n2)}
	假设 $f$ 是定义在圆周上的二次连续可微函数。当 $|n|\to\infty$ 时,有
	\[
	\widehat f(n)=O\left(\frac1{|n|^2}\right).
	\]
	则 $f$ 的 Fourier 级数绝对收敛且一致收敛于 $f$。
\end{corollary}

\begin{proof}
	当 $n\neq0$ 时,计算 Fourier 系数需要使用两次分部积分。取周期区间 $[0,2\pi]$,有
	\[
	\begin{aligned}
		2\pi\widehat f(n)
		&=\int_0^{2\pi}f(\theta)e^{-in\theta}\dd\theta                                      \\
		&=\left[f(\theta)\cdot \left(-\frac{e^{-in\theta}}{in}\right)\right]_0^{2\pi}
		+\frac1{in}\int_0^{2\pi}f'(\theta)e^{-in\theta}\dd\theta                              \\
		&=\frac1{in}\int_0^{2\pi}f'(\theta)e^{-in\theta}\dd\theta                              \\
		&=\frac1{in}\left[f'(\theta)\cdot\left(-\frac{e^{-in\theta}}{in}\right)\right]_0^{2\pi}
		+\frac1{(in)^2}\int_0^{2\pi}f''(\theta)e^{-in\theta}\dd\theta                          \\
		&=-\frac1{n^2}\int_0^{2\pi}f''(\theta)e^{-in\theta}\dd\theta.
	\end{aligned}
	\]
	由于 $f$ 和 $f'$ 是周期的,所以括号中的边界项均为零。因此
	\[
	2\pi |n|^2|\widehat f(n)|
	\leq \int_0^{2\pi}|f''(\theta)|\dd\theta
	\leq C,
	\]
	其中 $C$ 是与 $n$ 无关的常数。因为 $\sum 1/n^2$ 收敛,所以
	\[
	\sum_{n=-\infty}^{\infty}|\widehat f(n)|<\infty.
	\]
	由上一推论,$f$ 的 Fourier 级数绝对收敛且一致收敛于 $f$。
\end{proof}
进一步的,我们提高$f$的光滑性,得
\begin{corollary}\label{f(n)=O(1/n^k)}
	假设 $f\in C^k(\mathbb T)$,其中 $k\geq 1$。则当 $|n|\to\infty$ 时,
	\[
	\widehat f(n)=O\left(\frac1{|n|^k}\right).
	\]
	特别地,当 $k\geq 2$ 时,$f$ 的 Fourier 级数绝对收敛且一致收敛于 $f$。
\end{corollary}

我们还建立了以下重要等式：
\[
\widehat{f'}(n)=in\widehat f(n),
\qquad \forall n\in\Z.
\]
如果 $n\neq0$,这已经包含在上面的分部积分计算中；当 $n=0$ 时,
\[
\widehat{f'}(0)
=\frac1{2\pi}\int_0^{2\pi}f'(\theta)\dd\theta
=\frac{f(2\pi)-f(0)}{2\pi}=0=i0\widehat f(0),
\]
其中使用了 $f$ 的周期性。因此,如果 $f$ 是可微的且
\[
f\sim\sum a_ne^{in\theta},
\]
那么
\[
f'\sim\sum a_n in e^{in\theta}.
\]
同样,若 $f$ 是二次连续可微的,那么
\[
f''\sim\sum a_n(in)^2e^{in\theta}.
\]
在 $f$ 上附加更好的光滑性表现为 Fourier 系数衰减更迅速。
